% html2tex compatibility preamble (shared across ALL venue templates).
%
% This file collects every package load and macro provider that makes
% pandoc's HTML->LaTeX output compile cleanly under a venue's own style
% file. Each venue template (templates/<venue>/main.tex.template) should
% \input this file AFTER it loads its own \documentclass + venue .sty,
% and BEFORE \begin{document}. pack_tmlr_bundle.py copies this file into
% every bundle, so a new venue template inherits all accumulated fixes
% by adding a single \input line.
%
% Every entry below is load-bearing for a specific pandoc construct; the
% comment names the construct and the failure it prevents. Do NOT remove
% one without checking the matching post-processor rule in
% convert_to_tex.py.

% --- Captions / float numbering --------------------------------------
% The source HTML numbers its floats itself ("Table 8E", "Figure 5") in
% BOTH the prose and the caption text. LaTeX's automatic float counter
% counts differently (1..N in float order), which would desync caption
% numbers from the prose cross-references. Suppress the automatic
% "Table N:" / "Figure N:" label so the HTML's own number (kept in the
% caption text) is the single source of truth and always matches prose.
\usepackage{caption}
\captionsetup{labelformat=empty,labelsep=none,skip=4pt}

% --- Tables -----------------------------------------------------------
% tabularx + array: convert_to_tex.py rewrites pandoc longtable blocks to
% weighted tabularx so wide text wraps and columns are content-sized.
\usepackage{tabularx}
\usepackage{array}
% longtable: page-breakable tables. convert_to_tex.py emits tall
% verbatim-bearing tables (side-by-side review examples) as longtable so
% they split across pages instead of overrunning a single page / float.
\usepackage{longtable}
% booktabs: \toprule / \midrule / \bottomrule used by every converted table.
\usepackage{booktabs}

% --- Floats -----------------------------------------------------------
% float provides the [H] "put it exactly HERE, do not float" specifier.
% convert_to_tex.py places figures and (single-column) tables with [H] so
% they appear right where the source HTML put them -- right after their
% "Figure N shows..." sentence -- instead of floating off and (a)
% splitting a sentence across the float, (b) drifting into the
% references, or (c) leaving a near-blank page. In two-column mode a wide
% table still floats full-width ([t]) because [H] full-width is illegal.
\usepackage{float}
% placeins / \FloatBarrier kept available for any manual use, though the
% [H] policy means convert_to_tex.py no longer inserts barriers.
\usepackage{placeins}
% Vertical justification: \raggedbottom. tmlr.sty sets "\flushbottom \sloppy",
% but in a float-DENSE single-column paper \flushbottom is the wrong choice.
% Single-column tables are set [H] (exactly-here) and tall tables/figures
% FLOAT ([htbp]). When an [H] table or a deferred float forces a page break
% the page is left under-full, and \flushbottom stretches that leftover slack
% into a visible MID-PAGE band: a table at the top of the page, a big empty
% strip, then the following paragraphs pinned to the page bottom. That reads
% as a layout bug. \raggedbottom instead lets the leftover space collect at
% the page BOTTOM, so a float-forced short page looks like a normal short page
% (content from the top, clean whitespace below) rather than a torn one.
% Verified by render-and-measure over a 55-page float-dense build: switching
% to \raggedbottom removed every mid-page band (e.g. the strip after a mid-
% document results table) and introduced no gap-above-a-bottom-float
% regression. Prefer \raggedbottom for float-dense single-column venues;
% \flushbottom only suits text-dense papers whose pages fill naturally.
\raggedbottom

% Tighten the vertical gap around floats. Under the [H] policy a figure/table
% is set inline, so the space below its caption is \intextsep; LaTeX's ~12pt
% default reads as two blank lines after the caption. ~6pt is one clean line.
% \textfloatsep / \floatsep cover the full-width table* floats that still
% float to the column/page top. One change, every float and caption.
\setlength{\intextsep}{11pt plus 2pt minus 2pt}
\setlength{\textfloatsep}{12pt plus 2pt minus 2pt}
\setlength{\floatsep}{12pt plus 2pt minus 2pt}

% Table legibility (revision feedback): taller rows so consecutive entries are
% easy to tell apart, and more inter-column padding so adjacent cells never run
% together (e.g. a surname column touching the next column's text). One change,
% every table in the document.
\renewcommand{\arraystretch}{1.28}
\setlength{\tabcolsep}{8pt}

% Relax LaTeX's float-placement fractions. The defaults are conservative:
% \textfraction 0.2 forces >=20% of every float-bearing page to be text, so a
% large figure on a page with only a few lines of text is REJECTED and deferred
% -- the page is then left half empty (the canonical "big figure pushed to the
% next page, this page mostly blank"). \topfraction/\bottomfraction 0.7/0.3 cap
% how much of a page a top/bottom float may occupy, and \floatpagefraction 0.5
% sends a float to a near-empty dedicated page too eagerly. Relaxing all four
% lets a tall figure share a page with a little text (filling it) and lets a
% float sit at the page bottom, which is what removes the half-empty pages.
% Paired with [htbp] figure/tall-table placement (convert_to_tex.py).
\setcounter{topnumber}{4}
\setcounter{bottomnumber}{4}
\setcounter{totalnumber}{8}
\renewcommand{\topfraction}{0.92}
\renewcommand{\bottomfraction}{0.9}
\renewcommand{\textfraction}{0.06}
\renewcommand{\floatpagefraction}{0.6}
\renewcommand{\dbltopfraction}{0.92}
\renewcommand{\dblfloatpagefraction}{0.6}

% Never hyphenate cross-reference words: a broken "Fig-ure 2" / "Ta-ble 3"
% reads as an error, not a line break. A word listed with no discretionary
% hyphens is forbidden from breaking.
\hyphenation{Figure Figures Table Tables Appendix Appendices Section Sections}

% --- Graphics ---------------------------------------------------------
% graphicx: SVG figures are pre-rasterized to PDF (svglib) and referenced
% via \includegraphics{X.pdf}.
\usepackage{graphicx}
% xcolor: needed for the `black!75!blue` link-color mixing below (plain
% `color`, which hyperref loads, lacks the `!` mix syntax).
\usepackage{xcolor}
% svg: fallback only. We normally rewrite \includesvg{X.svg} ->
% \includegraphics{X.pdf}, but keep the package so a residual
% \includesvg (e.g. an SVG svglib could not convert) still compiles where
% Inkscape is available (Overleaf). Harmless when unused.
\usepackage{svg}
\svgsetup{inkscapelatex=false}

% --- Verbatim ---------------------------------------------------------
% fvextra extends fancyvrb with the `breaklines`/`breakanywhere` options
% that the base package lacks. convert_to_tex.py rewrites pandoc's plain
% \begin{verbatim} to \begin{Verbatim}[breaklines=true,breakanywhere=true]
% so long prompt/JSON/code lines WRAP instead of overflowing the page
% (the lmtt overfull-hbox warnings). Tables that contain a verbatim block
% are emitted as plain `tabular` (not tabularx, which double-scans its
% body and forbids verbatim).
\usepackage{fvextra}

% --- Math -------------------------------------------------------------
\usepackage{amsmath}
\usepackage{amssymb}

% --- Typographic quality ---------------------------------------------
% microtype: character protrusion + font expansion. Markedly improves
% justification, kills most overfull \hbox warnings, and is the single
% biggest "looks like a real paper" win. Guard the load so a venue that
% already loads it (or a non-pdftex engine) does not error.
\usepackage{microtype}
% Sensible link colors instead of hyperref's default garish boxes/reds.
% hyperref is loaded by the venue template BEFORE this file, so configure
% it here. Guard with \@ifpackageloaded so venues without hyperref skip.
\makeatletter
\@ifpackageloaded{hyperref}{%
  \hypersetup{%
    colorlinks=true,%
    linkcolor={black!75!blue},%
    citecolor={black!70!green},%
    urlcolor={black!60!blue},%
    breaklinks=true%
  }%
}{}
\makeatother

% --- Pandoc helper macros --------------------------------------------
% Pandoc defines these in its own default template; we ship our own
% template, so we provide them. \tightlist collapses spacing in lists
% pandoc marks "tight"; \passthrough wraps inline verbatim.
\providecommand{\tightlist}{\setlength{\itemsep}{0pt}\setlength{\parskip}{0pt}}
\providecommand{\passthrough}[1]{#1}
% Typographic character macros pandoc emits (would otherwise need
% textcomp, which can clash with venue fonts). Provide inline.
\providecommand{\textquotesingle}{\char39}
\providecommand{\textasciigrave}{\char96}
% Pandoc longtable width helper (only appears if a longtable slips past
% the rewriter); make it a no-op multiplier so it never crashes.
\providecommand{\real}[1]{#1}

% --- Citations --------------------------------------------------------
% Citation STYLE is venue-specific and therefore NOT set here (this file is
% shared by every venue). A numeric venue (TMLR, IEEE) sets
% \setcitestyle{numbers,square,comma} in its own main.tex.template AFTER this
% \input, paired with convert_to_tex.py --citations numeric (numeric \bibitem
% + [N] links). An author-year venue (TACL/ACL) sets nothing -- natbib's
% author-year default stands -- paired with --citations author-year
% (\bibitem[Author(Year)]{key} + \citep/\citet), which those venues REQUIRE
% (numeric citations are a desk-reject violation there).
